% =========================================================================
% MASTER BIBLIOGRAPHY: 30 VERIFIED ROCK-SOLID REFERENCES (MDPI / CATEGORY W)
% Strictly ascending citation keys: ref1 to ref30
% =========================================================================
\begin{thebibliography}{99}

\bibitem{ref1}
Cleland-Huang, J., Settimi, R., Zou, X., \& Solc, P. Automated classification of non-functional requirements. \emph{Requirements Engineering} \textbf{2007}. \url{https://doi.org/10.1007/s00766-007-0045-1}

\bibitem{ref2}
Kurtanović, Z., \& Maalej, W. Automatically Classifying Functional and Non-functional Requirements Using Supervised Machine Learning. \emph{2017 IEEE 25th International Requirements Engineering Conference (RE)} \textbf{2017}. \url{https://doi.org/10.1109/re.2017.82}

\bibitem{ref3}
Glinz, M. On Non-Functional Requirements. \emph{15th IEEE International Requirements Engineering Conference (RE 2007)} \textbf{2007}. \url{https://doi.org/10.1109/re.2007.45}

\bibitem{ref4}
Casamayor, A., Godoy, D., \& Campo, M. Automated Categorization of Non-Functional Requirements Using Semi-Supervised Learning. \emph{Information and Software Technology} \textbf{2010}. \url{https://doi.org/10.1016/j.infsof.2009.11.002}

\bibitem{ref5}
Zhao, L., Alhoshan, W., Ferrari, A., Letsholo, K. J., Ajagbe, M. A., Chioasca, E. V., \& Batista-Navarro, R. T. Natural Language Processing for Requirements Engineering: A Systematic Mapping Study. \emph{ACM Computing Surveys} \textbf{2021}. \url{https://doi.org/10.1145/3444689}

\bibitem{ref6}
Karlin, S., \& Burge, C. Dinucleotide relative abundance extremes: a genomic signature. \emph{Trends in Genetics} \textbf{1995}. \url{https://doi.org/10.1016/s0168-9525(00)89076-9}

\bibitem{ref7}
Shannon, C. E. A Mathematical Theory of Communication. \emph{The Bell System Technical Journal} \textbf{1948}. \url{https://doi.org/10.1002/j.1538-7305.1948.tb01338.x}

\bibitem{ref8}
Mantegna, R. N., Buldyrev, S. V., Goldberger, A. L., Havlin, S., Peng, C. K., Simons, M., \& Stanley, H. E. Linguistic Features of Noncoding DNA Sequences. \emph{Physical Review Letters} \textbf{1994}. \url{https://doi.org/10.1103/physrevlett.73.3169}

\bibitem{ref9}
Blaisdell, B. E. A measure of the similarity of sets of sequences not requiring sequence alignment. \emph{Proceedings of the National Academy of Sciences} \textbf{1986}. \url{https://doi.org/10.1073/pnas.83.14.5155}

\bibitem{ref10}
Ghandi, M., Lee, D., Mohammad-Noori, M., \& Beer, M. A. Enhanced Regulatory Sequence Prediction Using Gapped k-mer Features. \emph{PLOS Computational Biology} \textbf{2014}. \url{https://doi.org/10.1371/journal.pcbi.1003711}

\bibitem{ref11}
Reimers, N., \& Gurevych, I. Sentence-BERT: Sentence Embeddings using Siamese BERT-Networks. \emph{Proceedings of the 2019 Conference on Empirical Methods in Natural Language Processing (EMNLP)} \textbf{2019}. \url{https://doi.org/10.18653/v1/d19-1410}

\bibitem{ref12}
Devlin, J., Chang, M. W., Lee, K., \& Toutanova, K. BERT: Pre-training of Deep Bidirectional Transformers for Language Understanding. \emph{Proceedings of NAACL-HLT 2019} \textbf{2019}. \url{https://doi.org/10.18653/v1/n19-1423}

\bibitem{ref13}
Vaswani, A., Shazeer, N., Parmar, N., Uszkoreit, J., Jones, L., Gomez, A. N., Kaiser, Ł., \& Polosukhin, I. Attention Is All You Need. \emph{Advances in Neural Information Processing Systems (NeurIPS 2017)} \textbf{2017}. \url{https://doi.org/10.48550/arxiv.1706.03762}

\bibitem{ref14}
Hey, T., Keim, J., Koziolek, A., \& Tichy, W. F. NoRBERT: Transfer Learning for Requirements Classification. \emph{2020 IEEE 28th International Requirements Engineering Conference (RE)} \textbf{2020}. \url{https://doi.org/10.1109/re48521.2020.00028}

\bibitem{ref15}
Winkler, J., \& Vogelsang, A. Automatic Classification of Requirements Based on Convolutional Neural Networks. \emph{2016 IEEE 24th International Requirements Engineering Conference Workshops (REW)} \textbf{2016}. \url{https://doi.org/10.1109/rew.2016.021}

\bibitem{ref16}
Chawla, N. V., Bowyer, K. W., Hall, L. O., \& Kegelmeyer, W. P. SMOTE: Synthetic Minority Over-sampling Technique. \emph{Journal of Artificial Intelligence Research} \textbf{2002}. \url{https://doi.org/10.1613/jair.953}

\bibitem{ref17}
He, H., \& Garcia, E. A. Learning from Imbalanced Data. \emph{IEEE Transactions on Knowledge and Data Engineering} \textbf{2009}. \url{https://doi.org/10.1109/tkde.2008.239}

\bibitem{ref18}
Tantithamthavorn, C., Hassan, A. E., \& Matsumoto, K. The Impact of Class Rebalancing Techniques on the Performance and Interpretation of Defect Prediction Models. \emph{IEEE Transactions on Software Engineering} \textbf{2020}. \url{https://doi.org/10.1109/tse.2018.2876537}

\bibitem{ref19}
Agrawal, A., \& Menzies, T. Is 'Better Data' Better Than 'Better Data Miners'?: On the Benefits of Tuning SMOTE for Defect Prediction. \emph{Proceedings of the 40th International Conference on Software Engineering (ICSE)} \textbf{2018}. \url{https://doi.org/10.1145/3180155.3180197}

\bibitem{ref20}
Japkowicz, N., \& Stephen, S. The Class Imbalance Problem: A Systematic Study. \emph{Intelligent Data Analysis} \textbf{2002}. \url{https://doi.org/10.3233/ida-2002-6504}

\bibitem{ref21}
Lundberg, S. M., \& Lee, S. I. A Unified Approach to Interpreting Model Predictions. \emph{Advances in Neural Information Processing Systems (NeurIPS 2017)} \textbf{2017}. \url{https://doi.org/10.48550/arxiv.1705.07874}

\bibitem{ref22}
Ribeiro, M. T., Singh, S., \& Guestrin, C. "Why Should I Trust You?": Explaining the Predictions of Any Classifier. \emph{Proceedings of the 22nd ACM SIGKDD International Conference on Knowledge Discovery and Data Mining} \textbf{2016}. \url{https://doi.org/10.1145/2939672.2939778}

\bibitem{ref23}
Arcuri, A., \& Briand, L. A Hitchhiker's Guide to Statistical Tests for Assessing Randomized Algorithms in Software Engineering. \emph{Software Testing, Verification and Reliability} \textbf{2012}. \url{https://doi.org/10.1002/stvr.1486}

\bibitem{ref24}
Demšar, J. Statistical Comparisons of Classifiers over Multiple Data Sets. \emph{Journal of Machine Learning Research} \textbf{2006}. \url{https://doi.org/10.5555/1248547.1248548}

\bibitem{ref25}
Kitchenham, B. A., Pfleeger, S. L., Pickard, L. M., Jones, P. W., Hoaglin, D. C., El Emam, K., \& Rosenberg, J. Preliminary Guidelines for Empirical Research in Software Engineering. \emph{IEEE Transactions on Software Engineering} \textbf{2002}. \url{https://doi.org/10.1109/tse.2002.1027796}

\bibitem{ref26}
Rahman et al. (2026) RequiBERT: Transfer learning for software requirements classifications. \emph{Expert Systems with Applications} \textbf{2026}. \url{https://doi.org/10.1016/j.eswa.2026.132726}

\bibitem{ref27}
Naseer et al. (2026) A Fused Back Translation and BERT Model for Optimizing Software Requirements Classification. \emph{Arabian Journal for Science and Engineering} \textbf{2026}. \url{https://doi.org/10.1007/s13369-026-11528-z}

\bibitem{ref28}
AlDhafer et al. (2026) Adaptive Data Augmentation Techniques for Software Requirements Classification Using Deep Learning. \emph{IEEE Access} \textbf{2026}. \url{https://doi.org/10.1109/access.2026.3673382}

\bibitem{ref29}
Limaylla-Lunarejo et al. (2025) Using Voting and Stacking Ensemble Techniques to Optimize Software Requirements Classification. \emph{ACM / IEEE International Symposium on Empirical Software Engineering and Measurement (ESEM)} \textbf{2025}. \url{https://doi.org/10.1109/esem64174.2025.00037}

\bibitem{ref30}
Ayed et al. (2025) Benchmarking BERT for Software Requirements Classification Across Diverse Datasets. \emph{IEEE/ACS 22nd International Conference on Computer Systems and Applications (AICCSA)} \textbf{2025}. \url{https://doi.org/10.1109/aiccsa66935.2025.11315456}

\end{thebibliography}